\documentclass[Total.tex]{subfiles}
\begin{document}
	\chapter{Introduction}
		\begin{itemize}
			\item The vibrating string;
			
			Solution of the following equation
			\begin{gather*}
				-kx(t)=m\ddot{x}(t)
			\end{gather*}
			
			This second order ordinary differentail equation becomes
			\begin{gather*}
				\ddot{x}(t)+c^2x(t)=0\qquad c=\sqrt{k/m}
			\end{gather*}
			
			The general solution is 
			\begin{gather*}
				x(t)=acos(ct)+bsin(ct)
			\end{gather*}
			
			Where $a,b$ are constant give by $x(0)$ and $\dot{x}(0)$.
			
			Furthermore,
			\begin{gather*}
				x(t)=\sqrt{a^2+b^2}cos(ct-\varphi)=Acos(ct-\varphi)
			\end{gather*}
			\item Standing and Traveling Waves;
			\begin{itemize}
				\item Standing Wave:
				\begin{gather*}
					内容...
				\end{gather*}
				\item Traveling Wave:
				\begin{gather*}
					内容...
				\end{gather*}
			\end{itemize}
			\item Harmonic and Superposition of Tones;
			\item 
		\end{itemize}
		\section{Solution to the Wave Equation}
		\section{Fourier Coeffiecient}
			Suppose $f$ is a \underline{Riemann integrable functon}
			\begin{gather*}
				f:T=\RR/2\pi \ZZ\rightarrow \CC\mbox{ or }\RR
			\end{gather*}
			We have
			\begin{gather*}
				\hat{f}(r)=\frac{1}{2\pi}\int_0^{2\pi} f(t)e^{-irt}dt=\frac{1}{2\pi}\int_T f(t)e^{-irt}dt
			\end{gather*}
			According to the research in 19th,
			
			\begin{theorem}
				(Fejer's Theorem)
			\end{theorem}
	\chapter{Fourier Series}
		\section{Fundamentals}
			\subsection{Fourier Coeffiecient}
				\begin{Def}
					Define:
					\begin{gather*}
						a_n:=\frac{1}{L}\int_0^L f(x)e^{-\frac{2\pi i}{L}x}dx;n\in \ZZ
					\end{gather*}
				\end{Def}
				\textbf{Remark} We now assume that $f$ is Riemann-integrable on $[0,L]$. 
			\subsection{Some of Concepts in Analysis}
				\subsubsection{Everywhere Continuous Functions}
				\subsubsection{Piecewise Continuous Functions}
			\subsection{Definitions}
				\begin{Def}
					Take the following notion of
					\begin{gather*}
						\hat{f}(n):=a_n
					\end{gather*}
					Then, the \textbf{Fourier series} is given formally by:
					\begin{gather*}
						\sum_{n=-\infty}^\infty \hat{f}(n)e^{\frac{2\pi inx}{L}}
					\end{gather*}
				\end{Def}
				Then, 
				\begin{gather*}
					f(x)\sim \sum_{n=-\infty}^\infty a_n e^{\frac{2\pi inx}{L}}
				\end{gather*}
				
				Especially: Suppose $f$ is integrable on $[-\pi,\pi]$, then the $n$th Fourier coefficient of $f$ is
				\begin{gather*}
					\hat{f}(n)=a_n=\frac{1}{2\pi}\int_{-\pi}^{\pi} f(\theta)e^{-in\theta} d\theta;n\in \ZZ
				\end{gather*}
				
				Thus
				\subsection{Calculating of Some Fourier Coefiecients}
					\subsubsection{Poisson Kernel}
						\begin{gather*}
							D_N(x)=\sum_{n=-N}^N e^{inx}
						\end{gather*}
						is called n\textbf{th} Dirchlet kernel. It can be seen as the Fourier series with property $a_n=\begin{cases}
							1&|n|\leq N\\
							0&elsewise
						\end{cases}$. 
						
						A \textbf{closed form formula} for Dirchlet kernel is
						\begin{gather*}
							D_N(x)=
						\end{gather*}
					\subsubsection{Dirichlet Kernel}
				\subsection{Uniqueness}
				\subsection{Convolutions}
			\section{Good Kernels}
			\section{title}
			\section{Convergence}
	\chapter{Fourier Transforms}
	\chapter{Finite Fourier Analysis}
\end{document}